\documentclass[11pt]{article}
\usepackage{amsmath,amssymb,amsthm}
\usepackage[margin=1in]{geometry}

\newcommand{\Q}{\mathbb Q}
\newcommand{\Z}{\mathbb Z}
\newcommand{\Zp}{\mathbb Z_p}
\newcommand{\LambdaIW}{\Lambda}
\newcommand{\XiIW}{\Xi_{\mathrm{IW}}}
\newcommand{\Char}{\operatorname{char}}
\newcommand{\Gal}{\operatorname{Gal}}
\newcommand{\CRE}{\mathrm{CRE}}
\newcommand{\CRCFT}{\mathrm{CRCFT}}

\title{Step 190: Iwasawa Main Conjecture Carrier Pivot}
\date{}

\begin{document}
\maketitle

\section*{Verdict}

\[
V_{\mathrm{iwasawa\_closes}}.
\]

The Iwasawa Main Conjecture carrier is a non-\(\CRE\) p-adic module-theoretic carrier.  Its parent residual closes natively:
\[
\XiIW=0
\]
by the Mazur--Wiles proof of the cyclotomic Iwasawa main conjecture over \(\Q\), in the standard componentwise scope.

\section*{Inherited Record Audit}

The RH-thread records do not contain a complete Iwasawa carrier package.  The inherited framework contribution is the Step 187 dichotomy: non-\(\CRE\) RH-analogous carriers with a proved native ledger close natively.  The BSD-side records contain p-adic transfer rows, including Skinner--Urban / Iwasawa main-conjecture entries, but those are BSD-context records rather than an RH-thread Iwasawa carrier.

\section*{Carrier Declaration}

Fix an odd prime \(p\).  Let
\[
\Q_\infty/\Q
\]
be the cyclotomic \(\Zp\)-extension and let
\[
\Gamma=\Gal(\Q_\infty/\Q)\simeq \Zp.
\]
For a topological generator \(\gamma\in\Gamma\), set \(T=\gamma-1\).  The Iwasawa algebra is
\[
\LambdaIW=\Zp[[\Gamma]]\simeq \Zp[[T]].
\]

Let \(Q_n\subset \Q_\infty\) be the \(n\)-th cyclotomic layer and let \(A_n\) be the relevant \(p\)-primary class-group component.  The native module is
\[
X_\infty=\varprojlim_n A_n,
\]
with transition maps given by norms.  After fixing a Dirichlet character component \(\chi\), write \(X_\infty^\chi\) for the corresponding \(\Lambda_\chi\)-torsion module.

The analytic native object is the Kubota--Leopoldt p-adic L-function
\[
L_p(\chi)\in \Lambda_\chi,
\]
defined up to a \(\Lambda_\chi^\times\)-unit and characterized by its p-adic interpolation of special Dirichlet \(L\)-values.

\section*{Parent Residual}

The characteristic ideal of the Iwasawa module is
\[
\Char_{\Lambda_\chi}(X_\infty^\chi).
\]
Define the typed Iwasawa residual by the divisor defect
\[
\XiIW^\chi
  =
  \operatorname{div}_{\Lambda_\chi}
  \bigl(\Char_{\Lambda_\chi}(X_\infty^\chi)\bigr)
  -
  \operatorname{div}_{\Lambda_\chi}
  \bigl(L_p(\chi)\bigr).
\]
Equivalently, \(\XiIW^\chi=0\) iff
\[
\Char_{\Lambda_\chi}(X_\infty^\chi)=(L_p(\chi))
\]
as ideals of \(\Lambda_\chi\), up to unit.

The native probes are finite-order characters of \(\Gamma\), height-one prime valuations of \(\Lambda_\chi\), and norm-compatible class-group components.  The dissolving probes are the p-adic special-value interpolations of \(L_p(\chi)\).

\section*{Native Closure Theorem}

\begin{theorem}[Iwasawa carrier native closure]
For the cyclotomic Iwasawa carrier over \(\Q\), in the classical character-component scope of the main conjecture,
\[
\XiIW=0.
\]
\end{theorem}

\begin{proof}
Mazur--Wiles, \emph{Class fields of abelian extensions of \(\Q\)} (Inventiones Mathematicae, 1984), proves the cyclotomic Iwasawa main conjecture over \(\Q\): the characteristic ideal of the relevant Iwasawa module equals the ideal generated by the Kubota--Leopoldt p-adic L-function.  This is exactly the zero condition in the definition of \(\XiIW\).  Therefore the framework residual closes natively.
\end{proof}

Wiles, \emph{The Iwasawa conjecture for totally real fields} (Annals of Mathematics, 1990), extends the main-conjecture framework to totally real fields.  Skinner--Urban, \emph{The Iwasawa Main Conjectures for \(GL_2\)} (Inventiones Mathematicae, 2014), supplies GL2 p-adic main-conjecture results under their stated hypotheses.  These are scoped extension records, not used to enlarge the cyclotomic conclusion beyond its declared carrier.

\section*{CRE and CRCFT Audit}

The carrier is not \(\CRE\) for Riemann's RH.  Closing \(\XiIW\) proves the p-adic Iwasawa main conjecture in the chosen setting.  It does not imply Riemann's complex RH, and it is not logically equivalent to the assertion that the nontrivial complex zeros of \(\zeta(s)\) lie on the critical line.

Therefore \(\CRCFT\) does not apply.  CRCFT is the foreclosure taxonomy for classically-RH-equivalent carriers.  The Iwasawa carrier is instead a non-\(\CRE\) carrier with an available native proof ledger.

\section*{Dichotomy Implication}

Step 190 supplies the third non-\(\CRE\) native closure data point:
\[
\text{Selberg/Maass},\quad
\text{Weil--Deligne},\quad
\text{Iwasawa Main Conjecture}.
\]
The framework dichotomy is strengthened: CRE carriers land in CRCFT foreclosure modes, while non-CRE carriers with proved native ledgers close in their own typed context.

\section*{Nonclaim Boundary}

This step does not prove Riemann RH.  It does not close \(\Xi_{BC}\).  It does not use p-adic interpolation to infer the complex zero-line statement.  It does not weaken the retained no-gos.  It imports the Mazur--Wiles native theorem as a closure certificate for the Iwasawa carrier only.

\end{document}

